% =============================================================================
% БИЛЕТ 5
% =============================================================================

\ticket{5}{}

\question{Непрерывность композиции непрерывных отображений. Теорема о промежуточных значениях для функции, непрерывной в области}

\begin{defn}[Композиция отображений]
    Пусть $X, Y, Z$ --- множества, $f : X \to Y$, $g : Y \to Z$ --- отображения. Тогда \textit{композицией} отображений $f$ и $g$ называется отображение $g \circ f : X \to Z$, определённое формулой
    \[
        (g \circ f)(x) = g(f(x)).
    \]
\end{defn}

\begin{theorem}[Критерий непрерывности]
    Пусть $(X, d)$ и $(Y, \rho)$ --- метрические пространства, $f : X \to Y$ --- отображение. Тогда следующие условия эквивалентны:
    \begin{enum}
        \item $f$ непрерывно на $X$;
        \item для всякого открытого множества $V \subseteq Y$ множество $f^{-1}(V)$ открыто в $X$.
    \end{enum}
    \begin{proof}
        \textbf{$1 \Rightarrow 2$.} Пусть $V$ открыто в $Y$ и $x \in f^{-1}(V)$. Тогда $f(x) \in V$, поэтому существует $\varepsilon > 0$ такое, что $B\bigl(f(x), \varepsilon\bigr) \subseteq V$. В силу непрерывности $f$ в точке $x$ найдётся $\delta > 0$, для которого $f\bigl(B(x, \delta)\bigr) \subseteq B\bigl(f(x), \varepsilon\bigr)$. Следовательно,
        \[
            B(x, \delta) \subseteq f^{-1}\!\bigl(B(f(x), \varepsilon)\bigr) \subseteq f^{-1}(V).
        \]
        Таким образом, $x$ --- внутренняя точка $f^{-1}(V)$, значит $f^{-1}(V)$ открыто.

        \textbf{$2 \Rightarrow 1$.} Зафиксируем произвольные $x \in X$ и $\varepsilon > 0$. Шар $B\bigl(f(x), \varepsilon\bigr)$ открыт в $Y$, поэтому по условию (2) его прообраз $f^{-1}\!\bigl(B(f(x), \varepsilon)\bigr)$ открыт в $X$ и содержит $x$. Следовательно, существует $\delta > 0$ такое, что $B(x, \delta) \subseteq f^{-1}\!\bigl(B(f(x), \varepsilon)\bigr)$. Тогда
        \[
            f\bigl(B(x, \delta)\bigr) \subseteq B\bigl(f(x), \varepsilon\bigr),
        \]
        что и означает непрерывность $f$ в точке $x$. В силу произвольности $x$, $f$ непрерывно на $X$. \qedhere
    \end{proof}
\end{theorem}

\begin{theorem}[Непрерывность композиции]
    Пусть $(X, d)$, $(Y, \rho)$, $(Z, r)$ --- метрические пространства, $f : X \to Y$ и $g : Y \to Z$ --- непрерывные отображения. Тогда композиция $g \circ f : X \to Z$ также непрерывна.
    \begin{proof}
        Возьмём произвольное открытое множество $W \subseteq Z$. Покажем, что $(g \circ f)^{-1}(W)$ открыто в $X$. Заметим:
        \[
            (g \circ f)^{-1}(W) = f^{-1}\bigl(g^{-1}(W)\bigr).
        \]
        Так как $g$ непрерывно, по критерию (теорема выше) $g^{-1}(W)$ открыто в $Y$. Далее, поскольку $f$ непрерывно, $f^{-1}\bigl(g^{-1}(W)\bigr)$ открыто в $X$. Следовательно, $(g \circ f)^{-1}(W)$ открыто в $X$ для любого открытого $W \subseteq Z$, что по тому же критерию означает непрерывность $g \circ f$.
    \end{proof}
\end{theorem}

\begin{theorem}[Больцано--Коши о промежуточных значениях (обобщение)]
    Пусть $(X, d)$ --- линейно связное метрическое пространство, $f : X \to \mathbb{R}$ --- непрерывная функция. Тогда для любых $a, b \in X$ и любого числа $C$, лежащего между $f(a)$ и $f(b)$, существует точка $c \in X$ такая, что $f(c) = C$.
    \begin{proof}
        Без ограничения общности считаем $f(a) < C < f(b)$ (случай $f(a) > C > f(b)$ рассматривается аналогично). В силу линейной связности $X$ существует непрерывный путь $g : [0,1] \to X$ с $g(0)=a$ и $g(1)=b$. Композиция $h = f \circ g : [0,1] \to \mathbb{R}$ непрерывна (как композиция непрерывных отображений). Для функции одной переменной $h$ на отрезке $[0,1]$ выполняется классическая теорема Больцано--Коши: так как $h(0)=f(a) < C < f(b)=h(1)$, найдётся $t \in [0,1]$ с $h(t)=C$. Тогда $f(g(t)) = C$, и искомая точка --- $c = g(t)$.
    \end{proof}
\end{theorem}


\question{Разложение функции $(1 + x)^{\alpha}$ в ряд Тейлора}

% Ответ...

